\section{Registro crítico de cambios realizados}
\label{sec:evolucion}
Este registro documenta la sustitución del estado inicial A/B/C por el flujo
actual. Los ocho números corresponden a los dos bloques de implementación
acordados. Los puntos E01--E12 conservan el diagnóstico adicional. Su ejecución y la de
los antiguos puntos 9--14 se organiza ahora mediante P01--P20 en la
sección~\ref{sec:plan}; no son tres listas de trabajo independientes.

\subsection{Cómo se juzga una mejora}
Se distinguen cuatro tipos de conclusión:
\begin{description}
\item[Corrección comprobada:] una prueba reproduce una conducta incorrecta y
comprueba la propiedad esperada en el código actual.
\item[Optimización comprobada por conteo:] se comprueba que se realizan menos
cargas o llamadas equivalentes. Esto no acredita un porcentaje de reducción de
latencia, memoria total o coste real.
\item[Capacidad añadida:] existe una ruta de ejecución antes ausente. Su utilidad
para el usuario necesita evaluación funcional y de relevancia.
\item[Decisión con contrapartidas:] mejora una propiedad y puede empeorar otra.
Debe registrarse la alternativa y su condición de aceptación.
\end{description}
La evidencia disponible comprende 32 pruebas aprobadas al finalizar el bloque 8,
con LangGraph real y proveedores/modelos simulados, y análisis sintáctico de 17
archivos activos y de pruebas. Se probaron concurrencia del juez e inicialización
de recursos; no se hizo una prueba de carga HTTP con modelos reales.

\subsection{Cambio previo: nombre e identidad del proyecto}
Se sustituyeron las referencias textuales a la marca anterior por SPELLFIND en
12 archivos, incluidos interfaz, prompts, documentación y notebook; también se
renombró el archivo del grafo a \code{grafo_SPELLFIND.png}. Se conservaron los
identificadores de datasets externos y los títulos bibliográficos históricos.

La ventaja es la consistencia de nombres y referencias locales. No cambia el
algoritmo ni demuestra mayor calidad de recuperación. El renombrado del PNG no
implica que se haya redibujado o actualizado el contenido de la imagen.

\subsection{Punto 1: exclusiones persistentes}
\label{cambio:1}
\textbf{Antes.} A repetía la búsqueda final con un conjunto vacío de exclusiones,
de modo que una canción rechazada podía reaparecer. C añadía también las canciones
aprobadas a sus exclusiones.

\textbf{Ahora.} El grafo conserva un pool por solicitud y filtra el ranking al
seleccionar candidatos y al generar la respuesta. Solo los rechazos se agregan
al pool; el candidato aprobado queda primero. Si faltan candidatos, no se
reincorporan rechazados para completar cinco resultados. Se conserva el motivo
de cada exclusión. \code{ExclusionPool.exclude} evita duplicar su registro cuando
el juez y el orquestador registran el mismo rechazo.

\textbf{Archivos.} \code{judge.py}, \code{workflow.py} y \code{app.py}.
\textbf{Evidencia.} Las pruebas de rechazos integrados y de agotamiento verifican
que los IDs excluidos no vuelven, incluso después de refinar consultas.

\textbf{Balance.} Corrección de coherencia comprobada. La política conserva
también rechazos posiblemente equivocados del LLM; aumentar su precisión es
distinto de hacer cumplir sus decisiones. Un candidato rechazado bajo una
consulta anterior no se reconsidera después de refinar. E01 y E06 revisan esa
política y la calidad de la evidencia.

\subsection{Punto 2: veredictos explícitos y errores visibles}
\label{cambio:2}
\textbf{Antes.} A calculaba una variable de aprobación, pero devolvía verdadero
ante cualquier respuesta sin herramientas. Los errores del proveedor también
se convertían en aprobación.

\textbf{Ahora.} El juez devuelve \code{approved}, \code{rejected} o \code{error}.
Solo acepta aprobación textual explícita. Un rechazo textual reconocido o una
llamada válida a \code{exclude_song} rechaza al candidato. Respuestas vacías,
ambiguas, argumentos inválidos y fallos del proveedor producen error sin alterar
el pool. Todas las llamadas de herramienta se validan antes de excluir, incluido
que el ID coincida con el candidato actual.

Se eliminó la segunda llamada al LLM que solo confirmaba el rechazo: la exclusión
se resuelve después de validar la primera respuesta. La interfaz muestra una
advertencia cuando no se completa la validación. La API conserva el booleano
\code{approved} por compatibilidad y añade estado explícito y estado por resultado.

\textbf{Archivos.} \code{judge.py}, \code{app.py} y \code{templates/index.html}.
\textbf{Evidencia.} Pruebas de texto ambiguo, rechazo textual, herramientas
malformadas, ID ajeno y fallos de inicialización/proveedor.

\textbf{Balance.} Se corrigen falsas aprobaciones por errores de control. Esto no
demuestra que el LLM juzgue mejor una letra. El parseo estricto puede convertir
una respuesta semánticamente válida pero distinta del formato en error; no hay
reintento limitado ni validación humana del juez. Pendiente E06.

\subsection{Punto 3: contrato único de planes y fallback}
\label{cambio:3}
\textbf{Antes.} El prompt pedía listas, mientras que el fallback devolvía un
diccionario lírico incompatible con el worker.

\textbf{Ahora.} \code{search_contract.py} valida las cuatro listas. Adapta cadenas
individuales y el formato lírico antiguo con \code{text} y \code{weight}, elimina
espacios, textos vacíos y duplicados exactos, y rechaza elementos no textuales,
campos ausentes o planes sin consultas positivas. Los nulos se convierten en
listas vacías. El peso antiguo se valida pero no determina la fusión: los pesos
se siguen infiriendo en la aplicación. Router, fallback y refinador comparten
este contrato. El grafo puede usarlo incluso si falla la importación del router.

\textbf{Archivos.} \code{search_contract.py}, \code{router.py},
\code{workflow.py} y \code{refiner.py}.
\textbf{Evidencia.} Pruebas de esquemas antiguos, nulos, listas mixtas, duplicados,
falta de positivas y llegada del fallback a los workers y al juez.

\textbf{Balance.} Se comprueba compatibilidad de tipos, no calidad semántica.
El fallback copia la consulta original en ambas modalidades, sin traducir al
inglés ni preservar negativos detectados por un plan fallido. Puede empeorar
la recuperación y activar una modalidad innecesaria. Pendiente E02.

\subsection{Punto 4: aislamiento por solicitud}
\label{cambio:4}
\textbf{Antes.} A guardaba el pool activo en \code{_active_pool}; dos solicitudes
intercaladas podían modificar el pool de la otra.

\textbf{Ahora.} Se elimina esa variable. El pool se pasa explícitamente al juez;
la herramienta describe el rechazo y la aplicación lo aplica al contexto local
tras validarlo. Cada búsqueda tiene \code{request_id}, logs y copias al consultar
el conjunto de exclusiones. Una búsqueda nueva inicia sin rechazos previos.

\textbf{Archivos.} \code{judge.py}, \code{workflow.py} y \code{app.py}.
\textbf{Evidencia.} Una prueba con barrera intercala dos evaluaciones y verifica
que cada pool contiene solo su propio ID; otra ejecuta búsquedas sucesivas.

\textbf{Balance.} Elimina el cruce de exclusiones reproducido. No certifica que
todos los clientes externos o modelos soporten cualquier nivel de concurrencia.
Además, el pool se muta desde el juez y desde el grafo; la idempotencia evita
duplicados, pero conviene asignar un único propietario del estado. Pendiente E08.

\subsection{Punto 5: consolidación en un único grafo}
\label{cambio:5}
\textbf{Antes.} La web híbrida y el grafo semántico tenían reglas diferentes.
B era un borrador con incompatibilidades y C no recuperaba audio.

\textbf{Ahora.} \code{workflow.py} contiene el flujo común con estado tipado y
dependencias inyectadas: planificar, recuperar, evaluar y refinar. Web y consola
lo invocan. \code{cli.py} acepta una consulta por argumento. Los nombres antiguos
\code{tarea_3.py} y \code{tarea_3_correcion.py} redirigen a esa CLI y requieren
también un argumento; ya no ejecutan las consultas de ejemplo embebidas.
Los originales se conservaron en \code{experiments/}, marcados como históricos.

\textbf{Ventaja.} Hay una sola definición de exclusión, límite, refinamiento y
error para las entradas soportadas. La inyección permite ejecutar el grafo real
con modelos simulados en pruebas. Se verifican por inspección las entradas
compartidas; no se ejecutó una sesión CLI completa con modelos reales.

\textbf{Contrapartida.} LangGraph pasa a ser dependencia del flujo web; el grafo
se compila de nuevo por solicitud. Frente a C, el flujo nuevo no conserva
\code{MemorySaver}, streaming de actualizaciones por CLI ni recuperación de un
checkpoint tras una excepción. Por ello es una consolidación de lógica, no un
superconjunto de todas las capacidades previas. Pendiente E04.

\subsection{Punto 6: refinamiento selectivo por modalidad}
\label{cambio:6}
\textbf{Antes.} A descartaba candidatos con consultas fijas. C reescribía solo
letras y acumulaba vectores, con riesgo de repetir o arrastrar consultas fallidas.

\textbf{Ahora.} Después de un rechazo, \code{refiner.py} recibe consulta original,
plan, candidato y motivo del juez. Propone un plan nuevo y las modalidades que
necesitan cambiar. La validación impide cambiar las otras listas, eliminar
negativos existentes, vaciar una modalidad activa o activar una inicialmente
ausente. La normalización elimina duplicados exactos. Se admiten cambios
líricos, acústicos o de ambos; un plan sin cambios también es válido.

El prompt pide conservar imágenes literales, reemplazar positivos en vez de
acumularlos e impedir restricciones inventadas. Estas son instrucciones al
modelo: el validador no puede certificar por sí solo que una frase nueva conserve
la intención. El plan inicial, el final y cada refinamiento quedan registrados.

\textbf{Archivos.} \code{refiner.py}, \code{workflow.py} y \code{app.py}.
\textbf{Evidencia.} Pruebas de cambios por modalidad, planes sin cambios,
conservación de negativas, cambios indebidos y fallo del refinador.

\textbf{Balance.} Capacidad añadida, con mejora de relevancia aún no demostrada.
Un rechazo lírico no demuestra un fallo acústico: el refinador declara que no
hay evidencia de audio. Puede reescribir descriptores a partir de la petición,
pero no juzgar cómo suena el candidato. También añade una llamada LLM después
de cada rechazo refinable, hasta dos con tres intentos. Si falla, hoy se detiene
la búsqueda aunque queden candidatos del plan anterior. Pendientes E01, E03,
E05 y E07.

\subsection{Punto 7: terminación con significado explícito}
\label{cambio:7}
\textbf{Antes.} C enviaba a la ruta denominada aceptación al alcanzar tres
veredictos incluso si todos eran insuficientes. A podía devolver resultados
sin una aprobación clara al terminar el ciclo.

\textbf{Ahora.} \code{approved} requiere un veredicto explícito;
\code{limit_reached} significa presupuesto agotado con candidatos restantes;
\code{no_candidates} significa que no quedan candidatos recuperados sin excluir;
\code{error} distingue fallos y añade la etapa. No se refina después del último
intento ni se excluye una canción aprobada. Si el ranking queda vacío, el juez
no se invoca. Una canción sin letra consume un intento y se excluye sin LLM.

\textbf{Archivos.} \code{workflow.py}, \code{app.py}, interfaz y CLI.
\textbf{Evidencia.} Pruebas de límite, agotamiento, errores y ausencia de letras.

\textbf{Balance.} Corrección de control comprobada. El presupuesto cuenta
selecciones de candidato, no solo llamadas al juez. La ausencia de letra no
debería equivaler siempre a fallo de una consulta acústica o instrumental.
\code{no_candidates} describe el ranking explorado, no demuestra que todo el
catálogo carezca de canciones válidas. Pendientes E01 y E10.

\subsection{Punto 8: reutilización de rankings, embeddings y recursos}
\label{cambio:8}
\textbf{Antes.} A repetía ambas búsquedas hasta tres rondas y una búsqueda final,
aunque las consultas fueran iguales. Servidor y workers cargaban copias de
índices, mapa y CSV; los negativos se recodificaban para cada positiva.

\textbf{Ahora.} Cada solicitud cachea resultados por modalidad, positivas,
negativas y top-K. Las exclusiones se aplican después, por lo que no invalidan
esa caché. Cambiar solo letras reutiliza audio y viceversa. El resultado final
usa el último ranking disponible sin otra recuperación.

\code{resources.py} comparte catálogo, mapa e índices por proceso, con
inicialización protegida. El lookup usa un diccionario por ID y devuelve una
copia de los metadatos. El catálogo valida IDs duplicados y normaliza valores
ausentes. Los workers conservan sus modelos bajo bloqueo y no leen el CSV.
Cada uno tiene una LRU de 512 embeddings, devuelve copias y protege la
codificación para impedir mutaciones compartidas y cargas concurrentes dobles.

\textbf{Evidencia.} Con tres evaluaciones y consultas sin cambios, las pruebas
comprueban una llamada a cada worker frente a las hasta cuatro de A. Con un
cambio de una modalidad, comprueban dos llamadas a esa rama y una a la otra.
Otras pruebas verifican carga única de recursos/modelo bajo concurrencia y
copias independientes de embeddings.

\textbf{Balance.} Menos trabajo redundante demostrado por conteo; no existe aún
un benchmark con modelos reales. Las caches consumen memoria y no sobreviven
al proceso. Se comparte por proceso, no entre procesos de servidor. Se mantienen
tanto DataFrame como diccionario de metadatos, por lo que no se ha eliminado
toda duplicación de representación. Los bloqueos serializan codificaciones
por modalidad y el catálogo no se actualiza en caliente. Pendientes E07 y E09.

\subsection{Cambios auxiliares y compatibilidad}
\begin{itemize}
\item La API añade \code{termination_reason}, \code{error_stage},
\code{initial_plan}, \code{refinement_log}, \code{workflow_trace} y
\code{validation_scope: lyrics}. Se mantienen los campos de resultados previos.
\item Se devuelven \code{n_lyrics} y \code{n_audio}, que la interfaz ya intentaba
mostrar. Son cantidades del último ranking por rama, no contadores acumulados
de operaciones o canciones únicas de todo el proceso.
\item El servicio valida consulta no vacía y k positivo. El endpoint aún llama
a métodos del cuerpo antes de validarlo completamente: no está resuelto el
manejo HTTP de cualquier tipo de JSON. Pendiente E11.
\item Las rutas de lectura de catálogo e índices se resuelven desde el módulo.
El router y el indexador todavía mantienen rutas relativas al directorio de
trabajo; no se declara independencia completa de ese directorio.
\item Los workers ahora aceptan conceptos de cinco caracteres o menos y filtran
solo vacíos. Evita perder términos como piano, pero su impacto en relevancia no
ha sido medido.
\item El flujo actual propaga fallos de recuperación y descarta un ranking
antiguo si falla una búsqueda tras refinar. Evita mostrar scores asociados al
plan incorrecto. Frente a A se pierde la devolución parcial si una sola rama
falla; esa política debe revisarse, pendiente E05.
\item Se añadieron README, documentación, pruebas y exclusión de temporales en
Git. No se reescribieron los datos del catálogo ni se midió una fórmula nueva
de fusión: el score porcentual sigue pendiente de corrección.
\end{itemize}

\section{Puntos extra por realizar o corregir}
\label{sec:extras}
Esta lista conserva el diagnóstico previo al bloque I. E10 tiene cubierta su validación numérica por P03 y mantiene presupuestos pendientes en P11; E11 tiene entorno y HTTP verificados en P01--P02, con validación integral pendiente en P20; E07 tiene protocolo y simulación en P04, con calidad real pendiente. Los demás puntos continúan pendientes. El orden de esta lista
es histórico; la prioridad de ejecución vigente es la del plan P01--P20.
Se registran alternativas
con una condición verificable, no como mejoras ya implementadas. Prioridad alta
significa riesgo de respuesta incorrecta o pérdida de disponibilidad; media
indica decisiones de rendimiento, operación o mantenibilidad por medir.

\subsection{E01. Separar validación lírica, acústica y de consultas instrumentales}
\textbf{Prioridad alta.} El juez sigue leyendo letras y el grafo excluye canciones
sin letra. Una consulta solo acústica puede necesitar un resultado instrumental
que hoy se descartaría; una aprobación no valida audio. Además, un rechazo
lírico no justifica por sí solo refinar CLAP.

\textbf{Acción.} Definir requisitos y evidencia por modalidad; emitir
\code{not_evaluated} cuando falte evidencia acústica, permitir rutas solo audio,
y seleccionar refinamientos a partir de feedback fundamentado. Evaluar la
reconsideración de rechazos dependientes de una consulta refinada manteniendo
fijas las restricciones explícitas del usuario.
\textbf{Aceptación.} Consultas instrumentales no se rechazan por ausencia de
letra y ninguna salida afirma validación acústica sin evidencia.

\subsection{E02. Fallback fiel al idioma, modalidades y restricciones}
\textbf{Prioridad alta.} El respaldo usa el texto original en ambas ramas; su
formato es correcto, pero no garantiza traducción, separación modal ni negativos.
\textbf{Acción.} Diseñar un fallback degradado que conserve el último plan válido
o restricciones verificables e informe qué interpretación no pudo hacerse.
No inventar una extracción fiable de negativos mediante reglas sin evaluarla.
\textbf{Aceptación.} Consultas con negaciones e idiomas distintos del inglés no
pierden silenciosamente sus restricciones cuando falla el router.

\subsection{E03. Proveniencia de negativos y activación de modalidades}
\textbf{Prioridad alta.} Conservar todos los negativos impide quitar uno
inventado por el LLM. Impedir activar modalidades ausentes bloquea la reparación
de un error inicial del router.
\textbf{Acción.} Separar restricciones explícitas del usuario de hipótesis del
modelo; bloquear las primeras y permitir revisar las segundas mediante evidencia.
\textbf{Aceptación.} Se puede reparar una clasificación modal incorrecta sin
eliminar restricciones explícitas ni añadir requisitos no pedidos.

\subsection{E04. Coste del grafo y funciones perdidas frente a C}
\textbf{Prioridad media.} El grafo se compila por solicitud y ya no hay checkpoints
ni streaming equivalentes a C. LangGraph no hace que la recuperación sea mejor
por sí mismo.
\textbf{Acción.} Medir el coste de compilación y valorar un grafo compilado
reutilizable con estado por solicitud. Si se necesita reanudar o transmitir
progreso, recuperar esas capacidades con contratos serializables. Comparar con
una máquina de estados simple si el grafo no aporta ventajas medibles.
\textbf{Aceptación.} Igualdad funcional de entradas, aislamiento concurrente y
medidas de coste; declarar explícitamente si se ofrece reanudación o streaming.

\subsection{E05. Recuperación parcial y fallos del refinador}
\textbf{Actualización: resuelto técnicamente por P09 en la
sección~\ref{sec:bloque3}. El diagnóstico siguiente corresponde al estado previo.}
\textbf{Prioridad alta.} Ahora un fallo de una rama detiene toda la recuperación;
un fallo del refinador impide evaluar el siguiente candidato ya disponible.
Son decisiones conservadoras, no una mejora automática de disponibilidad.
\textbf{Acción.} Diseñar un estado degradado por modalidad y una salida que
conserve evidencia del plan usado. Evaluar continuar con los rankings previos
cuando falla el refinador y reajustar pesos solo con una política validada.
\textbf{Aceptación.} Un fallo de CLAP no borra resultados líricos útiles sin
explicación y un candidato no se marca aprobado por omitir la rama fallida.

\subsection{E06. Calidad del juez, formato y extensión de la evidencia}
\textbf{Actualización: P10 implementa formato, citas, cobertura y reintento;
calibración humana pendiente. Véase~\ref{sec:bloque3} para el estado actual.}
\textbf{Prioridad alta.} El formato estricto evita aprobaciones accidentales,
pero no calibra el juicio. Solo se inspeccionan 3000 caracteres, se juzga el
top-1 y no hay reintento de formato. Los demás resultados quedan sin evaluar.
\textbf{Acción.} Comparar salida estructurada, reintento limitado y políticas
de evidencia por fragmentos con referencia humana. Mantener fallos como errores
y diferenciar abstención de rechazo. Evaluar rechazo persistente y revisión
después de refinar sobre ejemplos etiquetados.
\textbf{Aceptación.} Medir falsos rechazos y aprobaciones, cubrir pistas situadas
al final de la letra y no presentar todo el top-K como validado por evaluar uno.

\subsection{E07. Comprobar si refinar y cachear mejora el resultado global}
\textbf{Actualización: P12 tiene comparación simulada y política experimental;
la utilidad musical y las medidas reales siguen pendientes. Véase~\ref{sec:bloque3}.}
\textbf{Prioridad alta para afirmar superioridad.} Las pruebas verifican ramas y
conteos; no prueban más relevancia ni menor latencia de extremo a extremo.
El refinador añade llamadas y puede desviar la intención original.
\textbf{Acción.} Comparar el mismo baseline corregido con y sin refinamiento,
y medir latencia fría/caliente, llamadas, coste, memoria, relevancia y negativos.
Usar el mismo corpus, consultas y criterios; repetir y reportar dispersión.
\textbf{Aceptación.} Publicar una comparación reproducible y habilitar el
refinamiento solo donde su beneficio justifique coste y errores adicionales.

\subsection{E08. Un único propietario de exclusiones y estado serializable}
\textbf{Prioridad media.} Juez y grafo mutan el pool; la idempotencia evita
duplicados, pero el estado contiene objetos y caches que no se han validado para
checkpoints. \code{TypedDict} no impone validación en tiempo de ejecución.
\textbf{Acción.} Hacer que el juez devuelva decisiones y el orquestador sea el
único que aplique exclusiones. Separar estado serializable y recursos efímeros,
con invariantes de aprobación, rechazo y versión del plan.
\textbf{Aceptación.} Una exclusión tiene un solo punto de escritura; si se añade
persistencia, guardar y restaurar conserva las decisiones sin compartir pools.

\subsection{E09. Invalidación de caches, memoria y concurrencia real}
\textbf{Actualización: P16 instrumenta contención y cachés y prueba concurrencia
local con índices reales/modelos de control; invalidación y memoria real siguen
pendientes de P16/P17. Véase~\ref{sec:bloque4}.}
\textbf{Prioridad media.} Reiniciar es hoy la política de actualización. Índices,
modelos, DataFrame y diccionario consumen memoria por proceso; devolver un
DataFrame o mapa compartido permite que otro consumidor lo modifique.
\textbf{Acción.} Medir RSS, GPU y contención; valorar snapshots versionados,
lecturas inmutables, publicación atómica del catálogo, límites de caché medidos
y batch de embeddings. No eliminar bloqueos sin comprobar el modelo concreto.
\textbf{Aceptación.} Resultados coherentes con una versión del índice, consumo
acotado y carga concurrente representativa; nunca mezclar mapa e índice viejos
y nuevos.

\subsection{E10. Presupuesto de candidatos y validación numérica}
\textbf{Actualización: validación numérica cubierta por P03 y presupuestos
separados por P11; véanse~\ref{sec:bloque1} y~\ref{sec:bloque3}.}
\textbf{Prioridad alta.} El límite cuenta también candidatos sin letra y solo
se exploran rankings de tamaño fijo. Faltan controles numéricos completos de
normas, valores finitos y compatibilidad índice--embedding. La resta negativa
puede producir una dirección nula o poco informativa.
\textbf{Acción.} Separar intentos de evaluación, descartes de metadatos y
ampliación de recuperación; validar vectores y catálogo al cargar.
\textbf{Aceptación.} Un fallo de datos no consume todo el presupuesto útil;
una consulta degenerada no entra a FAISS como una búsqueda aparentemente válida.

\subsection{E11. Entorno reproducible y contrato HTTP completo}
\textbf{Prioridad alta.} No se ha comprobado una instalación limpia de todas las
dependencias fijadas. El entorno temporal de pruebas recibió instalaciones
superpuestas; la última salida de metadatos informó LangGraph 1.2.16 y
LangChain Core 1.4.3 aunque también se intentó instalar LangGraph 1.2.4. Esa
salida no certifica un lock limpio ni la matriz completa de compatibilidad.
\textbf{Acción.} Crear un entorno aislado reproducible y separar dependencias
directas y de desarrollo. Verificar imports reales, versiones y resolución;
probar HTTP con Flask real, JSON inválido, estados 400/500/503 e inicialización
WSGI. Completar rutas, configuración, guía de instalación y modo de despliegue.
\textbf{Aceptación.} Instalar y ejecutar desde cero con versiones inequívocas,
sin depender de la carpeta temporal de esta sesión.

\subsection{E12. Completar los puntos originales 9--14}
\textbf{Actualización: P13--P15 cubren puntuación, pesos, filtros y evidencia
visible, con límites de calidad explícitos. P16 registra un perfil local; carga
real e ingesta P19 siguen pendientes. Véase~\ref{sec:bloque4}.}
\textbf{Prioridad según alcance.} Siguen pendientes puntuación interpretable y
pesos ajustables, filtros exactos, evidencia por resultado, operación completa,
evaluación reproducible e ingesta desde YouTube. El grafo unificado no los da
por implementados. E06, E07 y E11 detallan trabajos relacionados, sin duplicar
la numeración acordada.
\textbf{Aceptación.} Cada bloque futuro identifica los puntos cubiertos y sus
pruebas, y actualiza este registro sin atribuir mejoras todavía no medidas.

\subsection{E13. Fidelidad semántica de citas y cobertura del respaldo}
\textbf{Prioridad alta. Estado pendiente tras el bloque II.} La procedencia
verifica una cita literal y su negación, pero no demuestra que el texto inglés
propuesto por un modelo tenga exactamente el mismo significado. Los controles
léxicos tampoco cubren todas las lenguas, negaciones o contrastes. Una hipótesis
mal traducida podría quedar inmovilizada si se etiqueta como explícita.
\textbf{Acción.} Evaluar equivalencia de restricciones con ejemplos anotados,
separar evidencia textual de confirmación semántica y medir cobertura y coste
de reformulación. Considerar entrada estructurada de restricciones cuando aporte
una decisión útil al usuario, sin certificar una traducción solo por citarla.
\textbf{Aceptación.} Medir pérdidas/adiciones de intención y tasa de abstención;
ampliar reglas o modelos solo con evidencia. Se integra en P10/P12, y en P15
si se decide exponer corrección de restricciones al usuario.

\subsection{E14. Versionado por cambios de orden sin cambio de intención}
\textbf{Prioridad alta. Diagnóstico histórico; corregido el 18 de octubre de 2026.}
La huella del plan depende del orden de sus listas. Un refinador que solo
invierte dos consultas crea nuevas versiones y libera exclusiones previas.
Con dos candidatos y rechazo constante se evaluó al primero tres veces:
IDs $[1,1,1]$, versión final 3 y límite agotado, sin evaluar el segundo.
\textbf{Acción previa al bloque III.} Canonizar la identidad de consultas y
restricciones equivalentes, separando cambios de presentación de cambios
efectivos. Conservar evidencia y restricciones explícitas.
\textbf{Aceptación.} Permutar elementos equivalentes no libera rechazos;
cambios efectivos justificados sí permiten reconsiderar candidatos.
Asignado al bloque II/P07, con regresiones en P11/P12. La corrección y prueba
de no repetición constan en la sección~\ref{sec:bloque3}.

\subsection{E15. Serialización sin validar en la ruta de aclaración}
\textbf{Prioridad media. Diagnóstico histórico; corregido el 18 de octubre de 2026.}
El retorno temprano de planificación que solicita aclaración conserva metadatos
sin comprobar su serialización finita. Una entrada controlada con un campo
numérico NaN alcanzó el estado final de aclaración; JSON estricto lo rechaza.
Esto contradice la garantía general de estado JSON de P05. No se observó una
respuesta real del proveedor con ese dato.
\textbf{Acción previa al bloque III.} Aplicar el mismo contrato a todas las
salidas o definir un esquema explícito de metadatos permitidos.
\textbf{Aceptación.} Probar planes listos y aclaraciones con números no finitos
y objetos no serializables; todos los estados aceptados admiten JSON estricto.
La evidencia de E14/E15 y hashes de fuentes están en
\texttt{validation/revision\_block2/report.json}; son controles locales con
LangGraph real y colaboradores simulados, sin medir relevancia musical.
Asignado al bloque II/P05/P08; ambas rutas ahora validan metadatos JSON finitos.
Las pruebas actuales constan en la sección~\ref{sec:bloque3}.

\subsection{E16. Registro de consola incompatible con CP1252}
\textbf{Prioridad media. Diagnóstico histórico; corregido el 18 de octubre de 2026.}
En esta consola Windows, la suite de 72 pruebas terminó con dos errores
\texttt{UnicodeEncodeError} al imprimir símbolos Unicode desde el router y
el worker lírico durante su inicialización. Ejecutar con
\texttt{python -X utf8 -m unittest discover -s tests} dio 72 pruebas aprobadas.
No se descargaron modelos: las inicializaciones empleaban colaboradores
simulados. El worker CLAP contiene un mensaje similar, identificado por lectura
pero no reproducido en esta ejecución.
\textbf{Acción previa al bloque III.} Usar registros compatibles con el destino
y documentar UTF-8 como alternativa operativa. Un mensaje no debe interrumpir
la inicialización. \textbf{Aceptación.} Verificar salidas CP1252 y UTF-8 sin
fallos causados por el registro. Complementa P01/P02; limita la afirmación de
portabilidad de su verificación sin invalidar los demás controles.
Asignado al bloque I/P01/P02. La continuación introduce salida compatible y
verifica las 96 pruebas actuales en ambas codificaciones; véase~\ref{sec:bloque3}.

\section{Mantenimiento de la documentación y trazabilidad}
\label{sec:mantenimiento}
Cada entrega futura debe actualizar el README operativo y este registro antes de
declararse terminada. La entrada de cada cambio debe incluir fecha, requisito,
archivos, comportamiento anterior y actual, motivo, pruebas ejecutadas,
limitaciones, compatibilidad y tareas adicionales. Los puntos extra deben
mantener su ID y registrar cuándo pasan de pendientes a implementados o
descartados, con evidencia.

La descripción de A/B/C permanece como comparación histórica. Las capacidades
actuales se identifican como flujo consolidado; no se deben corregir los
originales archivados para que parezca que ya tenían las mejoras. Los resultados
experimentales de las presentaciones conservan su fecha y contexto.

Una afirmación de mejor rendimiento requiere medidas; una de mejor calidad
requiere evaluaciones comparables. Las pruebas con modelos simulados acreditan
contratos y control, pero no sustituyen pruebas con Groq, Nomic y CLAP reales.
% Estado adicional conservando el diagnóstico histórico.
\paragraph{18 de octubre de 2026: bloque V.}
Se incorpora publicación transaccional por generaciones, contexto de snapshot
por consulta, cola persistente con propiedad temporal y adaptador opcional de
ingesta. Se evita mezclar recursos o publicar trabajos vencidos; la búsqueda sigue
operativa durante la ingesta. La decisión de no recuperar checkpoints de búsqueda
y los costes de copiar índices se documentan en la sección~\ref{sec:bloque5}.
La calidad real y la compatibilidad de vectores nuevos no están demostradas.
E17--E19 conservan retención, integración real e instrumentales como pendientes.

